\section{Prior Australian reference exposure}
An earlier Australian recruitment stage identified \AURefYears{} observation-supported candidate station-years from \AURefStations{} stations. The historical intake record states that its designated F2 formal test was closed; this audit does not run that model or calculate new F2 errors. Subsequently, a separate, earlier-method Australian evaluation used calendar-year cases from the same acquired source archive. All \AUFirstMatch{} of the candidate station-years appear in that case file; \AUFirstEvaluable{} had at least one case in the old evaluation's eligible case set, including all \AUFormerTestEvaluable{} of the \AUFormerTest{} originally designated F2 test-role station-years. The prior evaluation contained \AUFirstCases{} eligible masked cases across its wider roster. These are previously evaluated records for a different pipeline, not unopened target station-years for a new independent test.

A later Australian development analysis used July--June rather than calendar-year windows. Counting physical time overlap rather than matching the year labels alone, \AUHotOverlap{} of the original \AURefYears{} calendar station-years intersect a saved development window, of which \AUHotEvaluable{} intersect at least one eligible development case. A separate Arrow/Python recount reproduced these counts. Matching station and interval does not prove that masked days are identical, that the F2 model was evaluated, or that data leaked into source-model training. The original native-sampling criterion allowed 20 observed hours per day and does not establish a complete 24-hour physical reference. These previously opened Australian records cannot supply a fresh independent confirmation of the current pipeline.
